\documentclass[12pt]{article}
\usepackage[margin=0.8in]{geometry}
\usepackage{amsmath}
\usepackage{amssymb}
\usepackage{xcolor}
\usepackage{tikz}
\usepackage{enumitem}
\usepackage{makecell}

\title{\LARGE Q3(b)(ii): Corner Solution Explained\\When RHS $<$ LHS at $x = 100$}
\author{ECO 310}
\date{}

\setlength{\parskip}{0.5em}

\begin{document}

\maketitle

\section*{The Setup}

You have $\$100$ to invest between two assets:
\begin{itemize}
    \item Asset 1: invest $x$ dollars
    \item Asset 2: invest $(100-x)$ dollars
\end{itemize}

Three states of the world:
\begin{itemize}
    \item State A: probability $1/3$, returns $(r_1^A, r_2^A)$
    \item State B: probability $1/3$, returns $(r^B, r^B)$ (both assets same)
    \item State C: probability $1/3$, returns $(r_1^C, r_2^C)$
\end{itemize}

Expected utility:
$$EU(x) = \frac{1}{3}u(w_0 + r_1^A x + r_2^A(100-x)) + \frac{1}{3}u(w_0 + 100r^B) + \frac{1}{3}u(w_0 + r_1^C x + r_2^C(100-x))$$

\section*{First-Order Condition}

Taking the derivative with respect to $x$:

$$EU'(x) = \frac{1}{3}(r_1^A - r_2^A)u'(w^A) + \frac{1}{3}(r_1^C - r_2^C)u'(w^C)$$

where:
\begin{itemize}
    \item $w^A = w_0 + r_1^A x + r_2^A(100-x)$ (wealth in state A)
    \item $w^C = w_0 + r_1^C x + r_2^C(100-x)$ (wealth in state C)
\end{itemize}

Simplifying:
$$EU'(x) = \frac{1}{3}\left[(r_1^A - r_2^A)u'(w^A) + (r_1^C - r_2^C)u'(w^C)\right]$$

Multiply by 3 and rearrange:
$$EU'(x) = (r_1^A - r_2^A)u'(w^A) - (r_2^C - r_1^C)u'(w^C)$$

\newpage

\section*{Interior Optimum: Equality Holds}

At an \textbf{interior solution} where $0 < x^* < 100$, we set $EU'(x^*) = 0$:

$$(r_1^A - r_2^A)u'(w^A) = (r_2^C - r_1^C)u'(w^C)$$

Rearranging:
$$\frac{u'(w^C)}{u'(w^A)} = \frac{r_1^A - r_2^A}{r_2^C - r_1^C} \qquad \text{(*)}$$

\begin{center}
\colorbox{green!20}{
\begin{minipage}{0.9\textwidth}
\textbf{At an interior optimum: LHS = RHS}

This balances the marginal benefit of investing in asset 1 vs asset 2.
\end{minipage}
}
\end{center}

\section*{Corner Solution: Inequality Holds}

But what if equation (*) \textbf{cannot be satisfied} for any $x \in [0, 100]$?

\subsection*{Case 1: RHS $<$ LHS even at $x = 100$}

\begin{center}
\colorbox{red!20}{\Large
$\displaystyle\frac{r_1^A - r_2^A}{r_2^C - r_1^C} < \frac{u'(w^C)}{u'(w^A)} \quad \text{even when } x = 100$
}
\end{center}

\textbf{What does this mean?}

This inequality tells us that even at $x = 100$:
$$EU'(x) = (r_1^A - r_2^A)u'(w^A) - (r_2^C - r_1^C)u'(w^C) > 0$$

\textbf{Even at $x = 100$, the derivative is positive!}

\newpage

\section*{What Should the Decision-Maker Do?}

\begin{center}
\colorbox{blue!20}{
\begin{minipage}{0.9\textwidth}
\textbf{If $EU'(x) > 0$ even at $x = 100$:}

The decision-maker wants to \textbf{increase} $x$ even more.

\textbf{Optimal solution: $x^* = 100$}

Invest \textbf{everything} (\$100) in asset 1, invest \textbf{nothing} in asset 2.
\end{minipage}
}
\end{center}

\subsection*{Why?}

\textbf{Interpretation of $EU'(x) > 0$:}

The derivative tells us: ``If I increase $x$ by one dollar, my expected utility goes \textbf{up}.''

So:
\begin{itemize}
    \item Increasing $x$ (more in asset 1) $\Rightarrow$ EU increases $\Rightarrow$ good
    \item Decreasing $x$ (more in asset 2) $\Rightarrow$ EU decreases $\Rightarrow$ bad
\end{itemize}

Since $EU'(x) > 0$ \textbf{even at $x = 100$}, you want to keep going: maximize $x$.

The maximum possible value is $x = 100$.

\section*{Visual Explanation}

\begin{center}
\begin{tikzpicture}[scale=1.5]
% Axes
\draw[->] (0,0) -- (6,0) node[right] {$x$ (dollars in asset 1)};
\draw[->] (0,0) -- (0,4) node[above] {$EU(x)$};

% EU curve (increasing)
\draw[very thick, blue, domain=0:5.5] plot (\x, {1 + 0.4*\x});

% Mark x = 0
\node[below] at (0, 0) {$0$};
\filldraw[black] (0, 1) circle (2pt);

% Mark x = 100
\node[below] at (5, 0) {$100$};
\filldraw[red] (5, 3) circle (2pt) node[above right] {$x^* = 100$ (optimal)};

% Derivative arrows
\draw[->, thick, red] (1.5, 1.6) -- (2.5, 1.6) node[midway, above] {$EU'(x) > 0$};
\draw[->, thick, red] (3, 2.2) -- (4, 2.2) node[midway, above] {$EU'(x) > 0$};

% Label
\node[blue, right] at (5.5, 3.2) {$EU(x)$ is increasing};

\end{tikzpicture}
\end{center}

\textbf{Since $EU(x)$ is always increasing in $x$, the maximum is at the rightmost point: $x^* = 100$.}

\newpage

\section*{Summary of Cases}

\begin{center}
\begin{tabular}{|l|l|l|}
\hline
\textbf{Condition} & \textbf{FOC} & \textbf{Optimal $x^*$} \\
\hline
\makecell[l]{Equation (*) satisfied \\ for some $x \in (0, 100)$} & $EU'(x^*) = 0$ & \makecell[l]{Interior solution \\ $0 < x^* < 100$} \\
\hline
\makecell[l]{RHS $<$ LHS \\ even at $x = 100$} & $EU'(x) > 0$ for all $x$ & \makecell[l]{Corner solution \\ $x^* = 100$} \\
\hline
\makecell[l]{RHS $>$ LHS \\ even at $x = 0$} & $EU'(x) < 0$ for all $x$ & \makecell[l]{Corner solution \\ $x^* = 0$} \\
\hline
\end{tabular}
\end{center}

\section*{Economic Intuition}

\textbf{Why would RHS $<$ LHS at $x = 100$?}

This happens when \textbf{asset 1 is so much better than asset 2} that even when you've invested everything in asset 1 ($x = 100$), the marginal benefit of investing even more in asset 1 is still positive.

In other words:
\begin{itemize}
    \item Asset 1 has better returns or better risk characteristics
    \item The decision-maker would prefer to invest \textit{more than} \$100 in asset 1 if they could
    \item But they only have \$100 total, so the best they can do is $x = 100$ (all in asset 1)
\end{itemize}

\section*{Key Takeaways}

\begin{enumerate}[leftmargin=2cm, itemsep=1em]
    \item \textbf{Interior solution:} FOC gives equality, $EU'(x^*) = 0$

    \item \textbf{Corner solution:} FOC cannot be satisfied in $(0, 100)$, check boundaries

    \item \textbf{If RHS $<$ LHS at $x = 100$:} Then $EU'(x) > 0$ everywhere, so $x^* = 100$

    \item \textbf{If RHS $>$ LHS at $x = 0$:} Then $EU'(x) < 0$ everywhere, so $x^* = 0$

    \item \textbf{The derivative tells you which direction to go:}
    \begin{itemize}
        \item $EU'(x) > 0$ $\Rightarrow$ increase $x$ (want more of asset 1)
        \item $EU'(x) < 0$ $\Rightarrow$ decrease $x$ (want more of asset 2)
        \item $EU'(x) = 0$ $\Rightarrow$ you're at the optimum
    \end{itemize}
\end{enumerate}

\end{document}
